%  LaTeX support: latex@mdpi.com 
%  For support, please attach all files needed for compiling as well as the log file, and specify your operating system, LaTeX version, and LaTeX editor.

%=================================================================
\documentclass[minerals,article,submit,pdftex,moreauthors]{Definitions/mdpi} 

%=================================================================

% MDPI internal commands
\firstpage{1} 
\makeatletter 
\setcounter{page}{\@firstpage} 
\makeatother
\pubvolume{1}
\issuenum{1}
\articlenumber{0}
\pubyear{2023}
\copyrightyear{2023}
%\externaleditor{Academic Editor: Firstname Lastname}
\datereceived{22 February 2023} 
\daterevised{23 April 2023} 
\dateaccepted{} 
\datepublished{} 
\hreflink{https://doi.org/} % If needed use \linebreak
%\doinum{}

%=================================================================
% Add packages and commands here. The following packages are loaded in our class file: fontenc, inputenc, calc, indentfirst, fancyhdr, graphicx, epstopdf, lastpage, ifthen, lineno, float, amsmath, setspace, enumitem, mathpazo, booktabs, titlesec, etoolbox, tabto, xcolor, soul, multirow, microtype, tikz, totcount, changepage, attrib, upgreek, cleveref, amsthm, hyphenat, natbib, hyperref, footmisc, url, geometry, newfloat, caption

\graphicspath{{figures/}} % path to folder with figures
\usepackage{subcaption}
\usepackage{listings}
\usepackage{xcolor}
\usepackage{gensymb} % for \textdegree
\usepackage{tabularx}

\usepackage{listings}
\usepackage{xcolor}

\definecolor{deepblue}{rgb}{0,0,0.5}
\definecolor{deepred}{rgb}{0.6,0,0}
\definecolor{deepmagenta}{RGB}{195,12,214}
\definecolor{deepgreen}{RGB}{29,122,30}
\definecolor{mintgreen}{RGB}{192,249,192}
\definecolor{codepurple}{RGB}{204, 0, 153}
\definecolor{LightCyan}{rgb}{0.88,1,1}
\definecolor{GainsboroPurple}{RGB}{247,176,247}
\definecolor{LightSlateGrey}{RGB}{124,118,118}
%    
\lstdefinestyle{mystyle}{
    commentstyle=\color{LightSlateGrey},
    keywordstyle=\color{deepgreen}\bfseries,
    emphstyle=\color{deepred},
    numberstyle=\tiny\color{deepmagenta},
    stringstyle=\color{codepurple},
    basicstyle=\ttfamily\footnotesize,
    breakatwhitespace=false,         
    breaklines=true,                 
    captionpos=b,                    
    keepspaces=true,                 
    numbers=left,                    
    numbersep=5pt,                  
    showspaces=false,                
    showstringspaces=false,
    showtabs=false,                  
    tabsize=2
}
\lstset{style=mystyle}

%=================================================================
% Full title of the paper (Capitalized)
\Title{Coherence of Bangui magnetic anomaly with topographic and gravity contrasts across Central African Republic}

% MDPI internal command: Title for citation in the left column
\TitleCitation{Coherence of Bangui magnetic anomaly with topographic and gravity contrasts across Central African Republic }

% Author Orchid ID: enter ID or remove command
\newcommand{\orcidauthorA}{0000-0002-5759-1089} % Polina Add \orcidA{} behind the author's name
\newcommand{\orcidauthorB}{0000-0002-6461-1551} % Olivier Add \orcidB{} behind the author's name

% Authors, for the paper (add full first names)
\Author{Polina Lemenkova $^{1}$*\orcidA{} and Olivier Debeir $^{1}$\orcidB{}}

%\longauthorlist{yes}

% MDPI internal command: Authors, for metadata in PDF
\AuthorNames{Polina Lemenkova and Olivier Debeir}

% MDPI internal command: Authors, for citation in the left column
\AuthorCitation{Lemenkova, P.; Debeir, O.}
% If this is a Chicago style journal: Lastname, Firstname, Firstname Lastname, and Firstname Lastname.

% Affiliations / Addresses (Add [1] after \address if there is only one affiliation.)
\address{%
$^{1}$ \quad Laboratory of Image Synthesis and Analysis (LISA), École polytechnique de Bruxelles (Brussels Faculty of Engineering), Université Libre de Bruxelles (ULB). Building L, Campus du Solbosch, ULB -- LISA CP165/57, Avenue Franklin D. Roosevelt 50, Brussels 1050, Belgium.
}

% Contact information of the corresponding author
\corres{Correspondence: polina.lemenkova@ulb.be; Tel.: +32 471 86 04 59 (P.L.)}

% Current address and/or shared authorship
%\firstnote{Current address: Affiliation 3.} 

% Abstract (Do not insert blank lines, i.e. \\) 
\abstract{
The interactions between the geophysical processes and geodynamics of the lithosphere play a crucial role in the geologic structure of the Earth's crust. The Bangui magnetic anomaly is a notable feature in the lithospheric structure of the Central African Republic (CAR) resulted from a complex tectonic evolution. This study reports on the coherence in the geophysical data and magnetic anomaly field analysed from a series of maps. The data used here include raster grids on free-air altimetric gravity, magnetic EMAG2 maps, geoid EGM2008 model and topographic SRTM/ETOPO1 relief. The data were processed to analyse the correspondence between the geophysical and geologic setting in CAR region. Histogram equalization of the topographic grids was implemented by partition of the raster grids into the equal-area patches of data ranged by the segments with relative highs and lows of the relief. The original data were compared with the equalized, normalized and quadratic models. The scripts used for cartographic data processing are presented and commented. The consistency and equalization of topography, gravity and geoid data were based using GMT modules 'grdfft' and 'grdhisteq' modules. Using GMT scripts for mapping the geophysical and gravity data over CAR shows advanced approach to multi-source data visualization to reveal the relationships in the geophysical and topographic processes in central Africa. The results highlighted the correlation between the distribution of rocks with high magnetism in the central part of Bangui anomaly, and distribution of granites, greenstone belts, and metamorphosed basalts as rock exposure. The correspondence between the negative Bouguer anomaly (<-80 mGal), low geoid values (< -12 m) and the extent of the magnetic anomaly with the extreme negative values ranging from -1000 to -200 nT is identified. The integration of the multi-source data provides new insights into the analyses of crustal thicknesses and the average density of the Earth in CAR, as well as the magnitude of the magnetic fields with notable deviations caused by the magnetic flux density in Bangui area related to the distribution of mineral resources in CAR.}

% Keywords
\keyword{Programming; geophysics; GMT; mapping; Africa; geology; geodynamics; geomatics; data analysis; cartography} 

\PACS{91.25.Rt; 91.25.-r; 91.10.-v; 91.10.Da; 91.10.Fc; 91.10.Op; 91.25.Cw}
\MSC{86A30; 86A04; 86A60; 86A99; 86-08; 86-11; 76E20}
\JEL{Q00; Q01; O13; Q5; Q51; Q55; Q56}

%%%%%%%%%%%%%%%%%%%%%%%%%%%%%%%%%%%%%%%%%%
\begin{document}

\section{Introduction}

\subsection{Background}

The problem of feature matching in Earth studies can be referred as matching the extent, direction and intensity of the geophysical and geologic processes, objects and phenomena visualised on the maps. Earth point representing the particular observed geophysical situation and geologic features is associated with a coordinate position in the cartographic domain and a feature variable describing the appearance around this variable. For analysis of correlation and links between diverse geologic and geophysical variables, matched feature points and continued fields represented on the maps should maintain  similar regional appearance as well as relative spatial relationships with other processes, e.g., variation in topography or geoid and regional distribution of the geologic units. Analysis of correlation between geophysical and geological variables has extensive uses in integrated geophysical and seismic analysis \cite{NDAMNJIKAM2023104816}, hydrological and engineering geological studies \cite{AMPONSAH2022104558}, geophysical anomalies \cite{MATENDE2023104766}, mineral exploration \cite{STEUER2020104172}, or landslide hazard risk assessment in the areas with complex geology \cite{CALAMITA2023106984,DIMAIO2020105473}.

Each cartographic matching method supports either a specific data format, e.g., like shape files -- .shp in the ArcGIS \cite{LAR2023e01436,NGENE2022360,CHENG2022106052} or the tiled format for image processing in remote sensing software \cite{KONWEA2023e01597,ALKHERSAN2022104678} or a limited set of converted and imported data formats from the multi-source data \cite{SAHA2022104706,APEH2022104561,SHEBL2022765}. Scripting and programming methods also showed their effectiveness in matching tasks and coherence analysis when dealing with topographic and geophysical datasets since their optimise the workflow by smooth, automated and rapid approaches in data processing. For instance, scripts facilitate modelling the geochemical-geophysical inversion to investigate the issues of dynamic topography \cite{Jones}, enable automated and optimised isolines approximation and mesh gradation in topographic data processing \cite{GORMAN20081721}, or support detailed topographic analysis through 3D cross-sections \cite{Sobh,Lemenkova202301}.

In this paper, we address the problem of the unified cartographic matching framework that supports a large family of data formats technically implemented in the Generic Mapping Tools (GMT) and QGIS. Our goal is to present a scripting cartographic approach that exploits the structural regularities in the geophysical, geologic and topographic datasets on Central African Republic (CAR) and uses them as input information during matching procedure. With this regard, this study aims at presenting a cartographic interpretation of the Bangui Magnetic Anomaly and analysis of its consistence using the available data on geophysical and geological setting of the CAR. Technically, we optimised the operational workflow of the cartographic work due to the use of scripts. 

To this end, we introduced the scripting cartographic framework that utilises the codes written by the GMT syntax for automation of the cartographic workflow. We first extract all the information from the gravity, topography and magnetic grids and map their parameters. Afterwards, we detect and separately analyse the repeated structures in the topographic grids. This is done using equalisation and comparison of the equalised, quadratic, normalized and original rasters through enhanced  ranges of a grid depending on frequency distribution of elevation amplitude of the relief in various regions of CAR. The raster images were incorporated independently and processed as separate files. The scripts used for implementation of the cartographic work are presented in the Appendix section with added comments on the most essential details of the GMT codes in Methodology.

Moreover, the coherency between the vertical gradient and the distribution of gravity anomalies is compared using several plotted map components. We used the data that are correctly matched and visualised for analysis of the repeated structures placed on the images and identified for the maps of gravity, geoid, topography, geology and magnetic anomalies over CAR. The main application for such an algorithm is for image pairwise comparison in map pairs where the overlap between the geological and geophysical parameters is comprised mostly of similarities in patterns on the analysed images. Our cartographic framework exploits five important characteristics of magnetic, geophysical and topographic grids along with geological data aimed at visualising the repeated structures, depressions and extremities in the geophysical anomalies over CAR. 

We demonstrate our cartographic framework performance in data processing works well for evaluating the coherency and feature similarity in geophysical, topographic and geological data over CAR. Our hypothesis is that a correlation exist between the geophysical and topographic processes that are tightly connected with regional geologic structures. Such relationship is linked to the tectonic processes that are expressed in the present land features and might be traced in the topographic surface and gravity variations. The effects of the geophysical anomalies on the distribution of the dependent variables can be observed as regional geomorphic structures. We showed that our approach performs effectively for the analysis of the coherency and revealing the links between the geophysical variables in the Earth critical zone in central Africa. 

\subsection{Study Area}

The Bangui magnetic anomaly is located in the Central African Republic (CAR) in the West Zaire Precambrian Belt (Figure \ref{fig01}). With an amplitude of 28 nT and the area of $700,000 km^2$ \cite{girdler1992possible}, it is the largest in Africa and one of the most important magnetic anomalies on the planet. The Bangui anomaly is named after the capital of the Central Africa Republic (CAR), located at 4°22'24'' N and 18°33'46'' E, on the right bank of the Oubangui River \cite{Baticle,Villien1990}. The Oubangui River is a major tributary of the Congo River playing important role of the economy of CAR as an important transport artery \cite{Vennetier}. 

\begin{figure}[H]
\centering
	\includegraphics[width=14.0 cm]{Fig_01.jpg}
	\caption{Topographic map of CAR. Data source: GEBCO \cite{GEBCO}. Map source: authors.
	\label{fig01}}
\end{figure}

Revealed by the satellite magnetometer data and reported in previous studies \cite{Regan,OUABEGO201311,NJITEUTCHOUKEU2021104206}, the anomaly has a crustal origin and is strongly related to the tectonic evolution and Precambrian history of the Saharan metacraton and Congo Craton. Geologically, it is located in the region of the deposited highly metamorphosed Precambrian rocks and has a Precambrian basement with large masses of the magmatic rocks \cite{Kochemasov}. The formation of the anomaly is related to the lithosphere bending. \cite{NJITEUTCHOUKEU2021104299} reports the thick crust and a strong lithosphere beneath the anomaly, as well as high coherence between the Bouguer gravity and relief which correlates with the Moho depths near the anomaly.

At the same time, the comparison between the magnetic contrasts, geologic setting and geophysical data in the Earth's critical zones provides new insights into the formation of the continental margins and distribution of mineral resources. It furthermore allows a better interpretation of the effects of the past tectonic activity on the present geophysical processes \cite{FELIXTOTEU2002906}. Thus, regions earlier affected by the geodynamic processes are notable for negative magnetic anomalies, which suggests a strong correlation between the active tectonics and the origin of the magnetic anomalies which is associated with the lithosphere movements \cite{RAVAT199259}. For instance, \cite{ARKANIHAMED1985257} noted that the lithosphere beneath the sedimentary basins formed during the Gondwana break up has a lower geomagnetism compared to the regions under the cratons, which is related to the thermal demagnetisation of the high-temperature asthenosphere. Further, certain effects are noted from the magnetic mineralogy of the ascending magma in the lower crust.

The magmatism in Central Africa is a consequence of the continental rifting, opening of the Equatorial Atlantic and the Mesozoic-Cenozoic magmatic activity \cite{WILSON1992203}. The lithospheric extension and the associated ascent of magma resulted in the intrusion of the tholeitic basalts which have typical geochemical and isotopic signatures. The tectonic structure of CAR is notable for the distribution of the Central African Fold Belt which extends in the north-east direction from Cameroon through CAR towards Sudan and contributed to the tectonic magmatism. The Central African Fold Belt is formed as a result of the extension of the Congo Craton and associated collisions of the related crustal blocks in Cameroon and CAR \cite{TOTEU2022296}. The Mesozoic separation of African crustal blocks during the Early Cretaceous resulted in the formation of the Western Central African Rift System. As a result, the crude oil deposits are formed in the structured reservoirs of the fault blocks in the rift basins of CAR \cite{GENIK1992169}.

The anomalies in the lithospheric magnetic field present the data on the variation in the intensity of Earth's magnetisation which points at both explicit and hidden geophysical properties. In turn, data on magnetisation reveal the properties of the underlying rock, the tectonic history of Earth and help interpreting the key steps in the geologic evolution of the Earth \cite{Vervelidou}. At the same time, our understanding of these processes and the complexity of their interactions using Earth observation data can assist in practical decision making on natural resource management, sustainable environmental development and mineral exploration in CAR. 

\section{Regional Geology} 

The geology of CAR is notable by the presence of the two prominent greenstone belts formed during the Archaean Eon as  metamorphosed mafic volcanic sequences within the granite-gneiss volumes \cite{Poidevin1991LesCD}. The greenstone belts are located to the north of Bouca in the West Zaire Precambrian Belt (Figure \ref{fig02}) as narrow subparallel bands of the extrusive igneous rocks (basalts and andesites) placed on a sialitic basement \cite{Boulvert1983}. The first one is a 250-km long Bandas belt composed of volcanic and metasedimentary rocks \cite{POIDEVIN1981157} and the second is located in the west -- a 150-km long Bogoin-Boali belt with anomalously high gold deposits \cite{Dostal}. They further include the tholeiitic basalts, dolerites from the Proterozoic dyke swarms and sills and andesites formed during crystallization of the basaltic magma \cite{Poidevin1985}. 

The geochemistry of the greenstone belts also includes the metamorphic minerals of the greenschist facies, pyroxenes and olivines \cite{POIDEVIN199497}. The presence of back-arc tholeiites and arc-related greywackes argues for a compressive margin plate boundary for the greenstone belts pointing at high tectonic activity in the marginal areas of the Congolese Craton \cite{cornacchia1989existence}. Further, the arc-related ultramafic origin of the bulk rocks caused by complex tectonic activity in the past, is also proved by \cite{Tanko} who reported on the geochemistry and petrogenesis of ultramafic rocks in the Precambrian terrane of the northern CAR. They indicated the presence of olivine and pyroxene, magnesio-hornblende and magnetite in the West Zaire Precambrian Belt (Figure \ref{fig02}) pointing at intense Precambrian mafic magmatism processes in the past. The magnesian intrusive of Paleoproterozoic age is also found in the Tamkoro-Bossangoa Massif in the northwest intruded into a strike-slip shear zone. They are composed by the gneisses with grained quartz diorites and biotite granites \cite{jgg2014243}.

The deposits of the Bangui basin in the southern region of CAR are dominated by the granodiorites which belong to the foreland of the Pan-African Oubanguides belt in the stratum deposited during the Proterozoic period of the Precambrian (pCm) \cite{POIDEVIN1986581}. The metamorphic rocks of the Bossangoa-Bossembélé area in the north of the country consist of the sedimentary and igneous units which are indicative of granulite facies conditions deposited on an old Paleoproterozoic continental crust \cite{Tankoetal2021,Gerard}. The Precambrian carbonate platforms are widely distributed in CAR as the remained evidence of the paleoenvironment. For instance, these include the carbonate deposits is the Ombella-M'poko Formation consisted of the metamorphosed carbonate facies (calcite, dolomite and quartz) on top of the fluvioglacial sediments covered by siliciclastics \cite{TOYAMA2019108}. 

Later examples of the palaeogeographic formations include the Carnot -- a Mesozoic fluvio-lacustrine detrital formation located in the western part of CAR and presented by clastic material formed before the end of the Cretaceous \cite{CENSIER199947}. More details on alluvial deposits of Carnot sandstones is given by \cite{CENSIER1990385}. Active tectonic movements and volcanic eruptions during the Cretaceous resulted in distribution of the typical morphological bodies of the intrusives igneous rocks -- dykes and sills \cite{LOULE201388}.

\begin{figure}[H]
	\includegraphics[width=14.0 cm]{Fig_02.jpg}
	\caption{Geologic provinces in CAR. Data source: USGS. Map source: authors.
	\label{fig02}}
\end{figure}

Regional stratigraphy of CAR is illustrated in Figure \ref{fig03}. The stratigraphic succession is presented by the oldest rocks from the Precambrian (Pc) and following Paleozoic (Pz) periods; the Mesozoic successions comprised by rocks from the Cretaceous (K) and Lower Cretaceous (Kl) periods; Cenozoic (QT) successions comprised of rocks  from the Pleistocene (Qp), Holocene (Qe), Quaternary (Q) and Tertiary (T) periods. The Precambrian (pCm) metamorphic formation is most exposed  with dominating quartzite-schistose and upper quartzite (Bangui and Nola) series. The granitic bodies and associated gneisses of Lower Proterozoic age of Precambrian period intruding the upper series were tectonized during the period of orogeny and form the basement in the south-central CAR \cite{LAVREAU199069}. The greenstone belt of pCm, covered by a series of conglomerates, is distributed on a granitized gneiss basement and includes volcanites and schists \cite{CORNACCHIA1989221}. 

The southern region of CAR lithology continues the northern border of the Congo Craton and is presented by the Lower Cretaceous (Kl) and Precambrian (pCm) successions (Figure \ref{fig03}). It is mostly composed of tholeitic basalts and gabbros of mafic and ultramafic rocks corresponding to the oceanic basalts and gabbros associated to iron-rich sediments and gneisses \cite{TOPIEN2023104793}. The Lower Cretaceous (Kl) succession, the second widely distributed in CAR, includes the fluvio-lacustrine deposits on the Central African shield \cite{Boulvert1983}. Further, the interfluves of the dense fluvial network of rivers in CAR correspond to the series of plateaus with the presence of the polymorphic sandstones Tertiary (T) with the valleys composed of soft sandstone and argillites of Cretaceous (K) \cite{Mestraudetal1964}. Further, the Bangui region is characterised by the Plutonic massifs distributed in the Ngouaka-Gbago area which includes the gabbro, granodiorite, and granites experienced the processes of the greenschist facies metamorphism \cite{VONTO2020103831}.

\begin{figure}[H]
	\includegraphics[width=14.0 cm]{Fig_03.jpg}
	\caption{Geologic units and lithology in CAR. Data: USGS. Map source: authors.
	\label{fig03}}
\end{figure}

The Paleozoic (Pz) period is notable for the emergence and geomorphic erosion in the earlier reliefs which resulted in the upper hilly surface of argillites in the northern regions of CAR \cite{Boulvert1995}. The alluvium and colluvium from the Pleistocene (Qp) valleys with modern sediments of Holocene (Qe) form a basin of the modern Aouk river and its tributaries which create a natural border between CAR and Chad, Figure \ref{fig03}. The formations in the alluvial deposits also include notable intrusions of diamonds formed as a result of the ascent of the upper mantle and related magmatic processes during the geological history of CAR \cite{Censieretal1998,Censier,MALPELI2014249,berthoumieux1957mission}. The geological development with favourable environmental setting of CAR situated in the equatorial, humid tropical and subequatorial climate resulted in its diverse natural resources. Rich mineral resources include gold mining fields in the quartz veins in greenschist facies of the Paleoproterozoic formations \cite{KPEOU2020103825} and alluvial diamonds \cite{USGSreportFR2013}. The latter ones are distributed in the south-western and north-eastern regions in the Carnot and Mouka-Ouadda Sandstone formations \cite{MALPELI2014249,USGSreport}. 

The geomorphology of most of the CAR is presented by flat or hilly plateau with dominated savannah \cite{Nguimalet} and dense ombrophile forest in the southern regions of the country \cite{Suchel} and the extreme northeast regions of CAR are a steppe. The northern region near the Bamingui  and Bangoran rivers is presented by the sahel-tropical forest \cite{Spinage1988}. The distribution of crops and related agriculture activities correspond to the different types of soil within various regions of the country which is in turn largely controlled by the regional setting of the geologic basement \cite{Villien1988,Quantin1965,Prioul}. For instance, the organic matter in soil is controlled by the underlying geology as determining its chemical and physical nature. As a result, the extent and types of habitats and vegetation patterns vary accordingly \cite{CombeauQuantin,Quantin1962}. Further, carbonated bed formations favour the development of karsts with typical relief and limestone soils in the southern regions of the Oubangui Basin \cite{Salomon}.

%----------- section ----------->
\section{Materials and Methods}
This study is based on the several gravity and geologic databases: two gravity databases (EMAG2 and WDMAM), satellite-derived gravity data (EGM2008) and topography data: GEBCO/SRTM \cite{GEBCO,SRTM}, ETOPO1 \cite{NOAA} and IGPP \cite{SANDWELL20211059,Pavlis}. The geologic maps were built upon the USGS data using QGIS software version 3.22.1 \cite{QGIS_software}, while the other datasets were processed separately by the GMT scripts by codes presented below. The programming approaches used for mapping present an advanced alternative to the traditional cartography \cite{TRAORE2020103933,RUNGE2005311,TRAORE2021100114}. They operate with scripts which can either be used as a sequence for codes using diverse modules used for plotting the maps \cite{4814952,Lemenkova202230,8588718,Lemenkova202228}, or as instruments in remote sensing data processing \cite{9964623,Lemenkova202229,7577632,Lemenkova202227}.

The advantages of the presented GMT-based modular algorithms include a high level of automation in data processing. Thus, the GMT-based approach does not require generating the projects for data processing, as in traditional GIS, but use scripts and the syntax of GMT instead. This enables to perform data processing on the fly. Such an advantage enables to adjust only relatively little parameters in the multi-source target datasets mapping for adaptation, textual annotations, accuracy of projections and fine tuning in visualization. The rest is performed by the script automatically which collects the lines of GMT code and implements the workflow without much supervision avoiding the need of human intervention as demonstrated in the scripts presented in the Appendix. Such technical approach presents a significant benefit of GMT over the conventional software and makes it a valuable instrument for geological and geophysical mapping when several datasets should be processed in an integrated way for quick and accurate cartographic data processing.

%----------- subsection ----------->
\subsection{Histogram equalisation}

The equalisation of the topographic grids was determined by means of the 'grdhisteq' module of the GMT version 6.1.1 \cite{WesselSmith1991,Wessel} by four different approaches in modelling elevation data: original, equalised, quadratic and normalised. The ’grdhisteq’ module of the GMT models the data values and divides a grid file into the equal patches. The histogram equalisation was implemented using the topographic ETOPO1 image covering CAR with the four cases of the relief modelling. The topography was divided into the equal-area segments for approximation of the regional elevation setting using equalised, normalized, quadratic and original models. As a result, the smooth distribution of the topographic values was demonstrated using Gaussian approximation and revealed notable peaks and depressions in the elevation of CAR which are comparable with the geophysical and magnetic grids. To this end, a sequence of the different GMT auxiliary modules was used to plot the map: 'pscoast', 'makecpt', 'grdimage', 'psscale' as well as annotations added by a combination of 'echo' and 'pstext'. The parameters of the commands are defined for fine cartographic adjustments and presented as a full code as showed in Listing \ref{lst01}. The comparison of the processed topographic grids regularised the non-linear models in the hypsometry of CAR to evaluate the extremities in the data covering the regions and highlight the depressions and peaks in the elevation models of the country.

%----------- subsection ----------->
\subsection{Geoid}
The first processing step consisted of converting the ADF file into the GDR by 'grdconvert' using the tile to the extent of CAR applied to the area of 'n00e00' in the global gravity data (Listing \ref{lst02}). Afterwards, the extent of the data range was evaluated for the image using the 'gdalinfo' module and the Color Palette Table (cpt) was adjusted accordingly using the 'makecpt' module. Extracting the image was done using the 'grdimage' module with the extensions applied to the actual extent of the country (-R14/28/2/11.5), that is, considering a EGM2008 clipped and extended over the target area between 14°-28°N of parallels and 2°-11.5°E of meridian. All the isolines were added on the image and interpolated with a 0.5 m step and annotations every 1 m to visualise the trends in variations of the values. To obtain the cartographic grid, ticks and annotations of graticule, the 'psbasemap' module was used with relevant parameters adjusted by flags, as showed in Listing \ref{lst02}. The auxiliary elements were added using the 'pscoast' and 'pstext' modules for the map to be completed. The GMT script written for plotting the geoid and demonstrated in Listing \ref{lst02} was used as the basis for plotting the maps of topography, and vertical gravity gradient in CAR with adjusted elements in codes and changed datasets processed in scripts, accordingly.

%----------- subsection ----------->
\subsection{IGPP Earth Free-Air Anomaly}
Gravity data were extracted from the EGM2008 database with global coverage resulting from the spherical harmonic models of the Earth's gravitational potential by \cite{Pavlis}. Diverse data processing algorithms were applied using the GMT modules as showed below in Listing \ref{lst03}. The raster image was visualised using the 'grdimage' modules with isolines added as contours with 20 m interval by the 'grdcontour' module. The 'Haxby' colour palette was applied, which is specially designed for gravity grids and visualising the geophysical data. The colour palette was adjusted according to the amplitude and geometry of the observed values in the raster grids, which were checked using the 'gdalinfo' module. The remaining cartographic elements were added using the sequence of modules GMT such as 'psbasemap', 'pscoast', 'pstext' and more, as demonstrated in Listing \ref{lst03}.

\subsection{Coherency of geophysical grids}

The 'grdfft' module that couples different grids and performs the computation on the grids in a frequency domain using Fast Fourier Transform (FFT), the coherency between the vertical gradient and free-air gravity anomaly was plotted and visualised along with the grids on magnetic, topographic and gravity data. The comparison of several maps allows the analysis of the correspondence of the crustal structure in the CAR region, the estimation of the topographic heights and hypsometry, and the coherency between the gravity and vertical gravity gradient. Such data integration allows better understanding of the geophysical and topographic responses on the inner behaviour of the lithosphere and geologic structure of CAR region, and their correspondence with the magnetic contrasts. The analysis of the coherency used for plotting the graph was performed using the GMT module 'grdfft' which processes the data using the Fast Fourier Transform, estimates the cross-spectral for the co-registered gravity grids (mGal), and reports the result as functions of the symmetry reflection between the both grids. To this end, the data were compiled and processed using the code in Listing \ref{lst04}. 

%%%%%%%%%%%%%%%%%%%%%%%%%%%%%%%%%%%%%%%%%%
%\pagebreak
\section{Results}

In this section, we present the results of this study which we briefly summarize as follows: 1) maps of the geologic structure and units in the region of CAR which enabled to analyse the correlation between the geological and geophysical setting; 2) maps of the gravity and geoid variations for estimation of the contrasts; 3) coherence between the vertical gravity anomalies and gradient; 4) visualisation of the extent of the magnetic anomalies for analysis of the correspondence with the lithosphere structure of CAR based on the datasets on geologic provinces and units. 

The analysis of the topographic grid reveals a local depression in the equalised raster grid which corresponds to the geoid and gravity data in the central region of CAR. Furthermore, the data are contrasted in the equalised and quadratic models (Figure \ref{fig04}, b) and d)) compared to the original and normalised grids (Figure \ref{fig04}, a) and c)), where topographic derivatives reveal a saddle-shaped structure in the 1-min topographic ETOPO raster in the region of the magnetic anomaly. This proves the distribution of the ring-shaped formation around the region of magnetic anomaly that has an outer ring diameter of over 810 km and an inner ring diameter at 490 km, which is confirmed by \cite{GIRDLER199245}. 

\begin{figure}[H]
\makebox[1.0\linewidth][r]{%
%\centering
	\includegraphics[width=17.0 cm]{Fig_04.jpg}
}
	\caption{Histogram equalisation of the topography grids on CAR based on the ETOPO1 raster data. Data source: ETOPO1. Map source: authors.
	\label{fig04}}
\end{figure}

The reasons of the gravity anomalies are caused by the rugged surface, nonlinear and complex composition of the Earth's inner structure. Other factors include thinner lithospheric thickness and Moho depths associated with the isostatic surface topography   \cite{Stunff} which affects its gravitational field. Furthermore, the values of gravity are affected by the rotation of the Earth which results that values near the equator are higher compared to those on poles. 

Beneath the magnetic anomaly, the geoid in CAR has values varying in the range from -4 to -8 m, which is lower compared to the western region where the values above 16 m are recorded. Moreover, the change in the geoid undulations over the elevations compared to the topographic depressions might be affected related to the plate tectonics and the processes of cooling lithosphere, thermal conductivity and mantle mineralogy \cite{HAGER198397}. Thus, the geoid height are sensitive to the distribution of density in the Earth's mantle which explains the variations in the values with higher values in the south-western region of CAR and negative values in the eastern regions, respectively (Figure \ref{fig05}). Furthermore, the contrasts in the values of the geoid are influenced by the lithospheric heterogeneities \cite{HagerClayton} which pose constraints in the density material that are reflected in the surface variations of geoid \cite{RICHARDS2023117964}.

\begin{figure}[H]
\centering
	\includegraphics[width=14.0 cm]{Fig_05.jpg}
	\caption{Earth geoid map of Central African Republic. Data source: EGM2008 Global Earth Geoid \cite{Pavlis}. Magnetic declination (magnetic rose in lower right corner) is given for Votofo, CAR: +1.63° (east, positive). Map source: authors.
	\label{fig05}}
\end{figure}

The geoid grid was mapped and compared with the magnetic anomaly in Bangui area (Figure \ref{fig05}) to show the coherency with structural variations in the crust reflected in the gravity values of geoid. Thus, the distribution of the magnetic contrasts continue up to a depth of 160 km where crustal thickness is about 40 km \cite{DORBATH1985231}, indicating a close relation between the Bangui magnetic anomaly and of the geoid and geophysical data that result in similar patterns contrast due to the Panafrican orogenic events \cite{MOLOTOAKENGUEMBA2008214}. Since the Bangui magnetic anomaly continues within the Central African Orogenic Belt which is a subject of the cumulative effects of the multiple tectonic events during the geological development in Precambrian and later periods, it has a thin lithosphere structure beneath the anomaly that well correlates with the gravity data \cite{GOUSSINGALAMO2018326}. 

Gravity anomalies and amplitude contrasts (Figure \ref{fig06}) result from the lower lithospheric depth and crustal density which correspond to the heterogeneities in the mantle density \cite{Boukeke}. As a result, similar gravity pattern at the surface closely follow the tectonic structure in the Central African Orogenic Belt due to the active mantle flow and thermal convection \cite{Daly}. The analysis of the Figure \ref{fig06} corresponds to the earlier maps \cite{Albouy} and reveals the negative values from -80 to -60 mGal over the study area (Figure \ref{fig06}. In contrast, higher values over +50 mGal are clearly visible in the western and southwestern regions of CAR that correspond to the regional outcrops of the Lower Cretaceous rocks (Figure \ref{fig03}) in the Zaire geologic province (Figure \ref{fig02}).

\begin{figure}[H]
\centering
	\includegraphics[width=14.0 cm]{Fig_06.jpg}
	\caption{IGPP Earth Free-Air Anomaly on Central African Republic. Data source: EGM2008. Map source: authors.
	\label{fig06}}
\end{figure}

The geologic structure in this region includes the Paleozoic glacial clastic sedimentary formations which are formed, among others, by tillites, argillites and mudstones cropping out from several sedimentary formations \cite{MILESI2006571}. The comparative analysis of Figures \ref{fig06} and \ref{fig07} with data on magnetic anomaly shows that negative Bouguer anomaly of ca. -40 mGals (blue colours in Figure \ref{fig06} and dark blue areas in Figure \ref{fig07}) coincides with the distribution of the magnetic anomaly, which implies that these heterogeneities are related to the Earth lithosphere and geologic features increasing or depreciating local magnetic field depending on the subsurface structure and composition of the Earth's crust. Furthermore, the variations in vertical gravity intensity observed in Figure \ref{fig07} point at the difference in lithospheric mass anomalies caused by deeper density contrasts. 

Since surface topography is related to the geologic-tectonic processes, it is also associated with seismic anomalies in the dense mantle \cite{Chevrot}. In turn, seismic velocity anomalies are proportional to the heterogeneities in mantle density related to the tectonic activities \cite{Forte}. Therefore, the observed geoid heights over cratons differ in gravity values than in the region of the magnetic anomaly and the distribution of the related Central African Orogenic Belt. For instance, major structural elements in the Bossangoa and Bossembele areas in the northwest reflect the tectonic evolution of the basement of the North Equatorial fold belt which includes the shear zones, lineations and fold axes \cite{Mapoka} that affect the shape of the geoid and topography in the modern relief of CAR. As a result, gravity values affected by the crustal density and continental Iithospheric thickness are reflected in the topography and the geoid patterns of CAR. 

\begin{figure}[H]
\centering
	\includegraphics[width=14.0 cm]{Fig_07.jpg}
	\caption{Derivative of the IGPP Earth’s gravity field on Central African Republic: vertical gravity gradient. Data source: EGM2008 Global Earth Geoid \cite{SANDWELL20211059}. Map source: authors.
	\label{fig07}}
\end{figure}

The distribution of the bipolar magnetic anomaly has the major part of its values negative over the two third of CAR in the central and south-west regions of the country. It also showed a distinct large negative depression with values greater than 1000 nT and surrounding peaks which well correspond with previous records \cite{Taylor2007}. The negative values are visualised as dark blue colours (Figure \ref{fig08}) while the peaks are visible in the northern regions of the country with corresponding values of 300 to 400 nT (dark red colours in Figure \ref{fig08}). Moreover, these data conform the records from the Magsat, a satellite specially used for the measurement of the crustal magnetic anomalies of the Earth. It shows the existing correspondence between the magnetic data and the surface geology and tectonics in central Africa \cite{Green}. The magnetic intensity (Figure \ref{fig08}) compared with the Central African geologic and gravity maps (Figure \ref{fig02}, Figure \ref{fig03} and Figure \ref{fig04}) shows that most of the geologic units and geophysical patterns are reflected by the magnetic fields which proves the interrelation between the topography and the geologic structure with the magnetic field data over CAR. 

Lower magnetic values are observed across the gneissic and granitic rocks of the shields and across the greenstone belts of CAR. Furthermore, they correspond to the area over the greenstone belt intruded by the Precambrian rocks with ultrabasic dykes. In contrast, the extent of the East Africa Rift correlates with the moderate magnetic values which lie between the -20 to 10 nT (correspond to yellow colours in map of Figure \ref{fig08}) and from +20 to +40 nT in the north-east region of the country (soft orange colours in Figure \ref{fig08}) in Sud geologic Province with Pleistocene outcrops. The highest magnetic values over +60\hl{ }nT (bright red colours in Figure \ref{fig08}) correspond to the distribution of the Precambrian rocks in the south and occasional Tertiary rocks in the north \hl{of the country}. The most heterogeneous pattern of \hl{the }magnetic data is visible in the West \hl{Zaire} Precambrian Belt where the lowest depressions of \hl{the }dataset contrast with the highest peaks in values (Figure \ref{fig08} cf. Figure \ref{fig02}).

\begin{figure}[H]
\centering
	\includegraphics[width=14.0 cm]{Fig_08.jpg}
	\caption{Magnetic anomaly in CAR based on WDMAM with 3 arc minute resolution. Dark blue area indicate the negative Bangui magnetic anomaly in the centre of CAR.\hl{ }Data source: WDMAM \cite{Lesur}. \hl{Map} source: authors.
	\label{fig08}}
\end{figure}

Figure \ref{fig09} shows a detailed view of the anomaly derived from the 2-arc minute resolution grid. It reveals taht the Bangui anomaly consists in three separate segments with a large central negative anomaly having elongated shape and two with positive anomalies situated to north- and southwards of the large negative depression. The Bangui magnetic anomaly derived from the EMAG2 grid illustrates a typical interpolated pattern of the magnetic anomalous  field, magnetised by induction in a zero inclination field \cite{RAVAT2002131}. As the magnetic equator passes near the centre of the surface, the computed-based modelling reveals the distribution of the negative and positive values in Bangui anomaly accordingly. Certain trends exist in data on daily amplitude of the magnetic variation in Bangui magnetic anomaly which demonstrated fluctuations over the long-term period \cite{Zhang,Barbas}. These shown possible fluctuations in the Earth's magnetic field that might be affected by the ionospheric conductivities. However, the visualised dataset is based on the EMAG2 grid which is compiled from the remote sensing observations obtained from the satellite, ship, and airborne magnetic measurements with improved resolution to 2 arc minutes \cite{Maus}.
 
\begin{figure}[H]
\centering
	\includegraphics[width=14.0 cm]{Fig_09.jpg}
	\caption{Magnetic anomaly in CAR based on EMAG2 with 2 arc minute resolution. Dark blue area indicate the negative Bangui magnetic anomaly. Black and dark grey areas mean the "no data" areas where no magnetic survey was done in the original dataset.\hl{ }Data source: EMAG2 \cite{Meyer}. \hl{Map} source: authors. 
	\label{fig09}}
\end{figure}

Figure \ref{fig10} \hl{demonstrates} the coherency between the geophysical data, hypsometry and the Bangui magnetic anomaly fields. Thus, the gravity data with -120 mGal Bouguer anomaly, the associated negative depressions in the geoid derived from the EGM2008 and the ETOPO1-based topographic hypsometry over CAR all correlate with the negative values of the Bangui magnetic anomaly that are below -1000 nT. At the same time, the comparison of the extent of Bangui magnetic anomaly (Figures \ref{fig08} and \ref{fig09} with \hl{an }enlarged view\hl{ in }the composed plot \hl{shows} the coherency with geophysical data and the geologic maps of CAR (Figures \ref{fig03} and \ref{fig04}) demonstrates the predominant distribution of the Precambrian (pCm) rocks in this critical region that are exposed just under the central area of the magnetic anomaly and depressions in gravity. Spatial\hl{ }extent in the\hl{ }grids (Figure \ref{fig10}) \hl{showing} combined magnetic\hl{ }and Bouguer gravity data (WDMAM, EGM2008, IGPP and EMAG2) \hl{indicates} the integrity of the geophysical and geomagnetism in CAR.

\hl{ The} correlation is visible between the lowest magnetic data in the \hl{region from} 3°N \hl{to} 7°N (\hl{a }white square in Figure \ref{fig10}) and\hl{ the }elongated negative depression in the free-air gravity\hl{ }at 4°N over the basement of West \hl{Zaire} Precambrian Belt which\hl{ }has a south-west-north-east extension over the territory of CAR. Such correspondence is related to the variations in crustal thicknesses in this region of CAR caused by the tectonic processes\hl{. }This \hl{points at} the interrelation of gravity anomalies\hl{ }with amplitudes ranging from -85 to +80 mGal over CAR with negative values beneath the magnetic anomaly. \hl{Such} contrasts \hl{indicate} the extent of \hl{the }major tectonic structures of CAR:\hl{ }the Central African Orogenic Belt on the Precambrian shield and the northern part of the Congo Craton.

\begin{figure}[H]
\centering
	\includegraphics[width=13.5 cm]{Fig_10.jpg}
	\caption{Coherency between \hl{the }geophysical grids on CAR. \hl{\textbf{(a)} Coherency between the vertical gravity and free-air gravity anomaly grids on CAR; \textbf{(b)} Magnetic anomalies on CAR based on the EMAG2 grid; \textbf{(c)} Magnetic anomalies representing the WDMAM grid; \textbf{(d)} Vertical gravity gradient; \textbf{(e)} Free-air gravity anomalies over CAR; \textbf{(f)} Relief model based on the IGPP; \textbf{(g)} Geoid model based on the EGM2008. Map source}: authors. 
	\label{fig10}}
\end{figure}

%%%%%%%%%%%%%%%%%%%%%%%%%%%%%%%%%%%%%%%%%%
\section{Discussion}
Detecting the patterns in the lithospheric structure using the integrated mapping of gravity, magnetic fields, geoid and topography data is of paramount interest for \hl{the }Earth science. Accordingly, \hl{we examined the relationship between these datasets over CAR by the advanced cartographic approach. To this end, }we\hl{ }presented a GMT console-based method of \hl{the }cartographic data processing, which \hl{was} integrated into the geological-geophysical analysis of \hl{the Bangui magnetic anomaly and surrounding regions}. We\hl{ }tested our proposed techniques in the context of the \hl{regional} setting of CAR where extreme magnetic anomaly is presented in the Bangui area. Overall, seven datasets are processed, visualised, evaluated and analysed \hl{including the }geoid EGM2008, IGPP, \hl{GEBCO}, ETOPO1, EMAG2, WDMAM and \hl{the} geologic datasets \hl{obtained from the USGS}. 

\hl{The coherency between the geologic structure of CAR and geophysical datasets was determined by means of the multi-source datasets using diverse modules of GMT to assess the regions crossed by the Bangui magnetic anomaly. The results show that the gravity structures and the hypsometry associated with the Bangui magnetic anomaly are strongly linked with the extent of the major geologic units and tectonic structures including Central African Orogenic Belt and the northern part of the Congo Craton. }As \hl{a continuation of this }work \hl{in the future, three} directions are possible. First, to integrate other datasets \hl{including remote sensing data}. Second, to extend the current methods of cartographic scripting by other modules of GMT. Third, to \hl{extend the } study area \hl{towards the} neighbouring regions of Africa \hl{for a further geological analysis}. Furthermore, the estimated values of the magnetic lowest negative values between the -1000 nT to -280 nT were compared to the extent of the topographic and geophysical grids to reveal the correlations between the datasets. The obtained results \hl{are demonstrated on} the \hl{presented }maps \hl{which show }that\hl{ }gravity contrasts \hl{are }related to the \hl{extent of the }magnetic anomaly \hl{which are tightly connected to} the tectonic structures and geologic data on CAR. 

\hl{This study presented previously unpublished analysis of the coherence between the geophysical and geomagnetic data over CAR implemented by the GMT scripts. Limited information on the Bangui magnetic anomaly and scarce data on its relationship with regional geologic setting, }lithospheric structure and \hl{the gravity motivated this research and underlined the actuality of the presented results}. With this regard, mapping the magnetic and gravity anomalies \hl{over CAR }analysed in the context of \hl{its} regional geologic and tectonic development\hl{ }contributes both to the theoretical investigation on geotectonic and geological setting in \hl{central Africa} and to the practical increase of \hl{geospatial information} aimed at mineral exploration in CAR. With respect to \hl{the }practical geologic \hl{aims and management of }natural resources\hl{, }mapping the coherency between the Bangui magnetic anomaly, geophysical data and geologic setting\hl{ }contributes to the analysis of \hl{the }Earth's structure and\hl{ increase the knowledge over geophysical and geological setting in central Africa}.

%%%%%%%%%%%%%%%%%%%%%%%%%%%%%%%%%%%%%%%%%%
\section{Conclusions}

This \hl{study} describes the correlation between the distribution of rocks with high magnetism in the central part of the Bangui anomaly of CAR, geologic fields of granites, greenstone belts, and the metamorphosed basalts as rock exposures. Moreover, it presented and compared the data on \hl{the extent of the }negative Bouguer anomaly below the -80 mGal with low geoid values\hl{, as well as }the distribution of the magnetic contrast with \hl{abnormally} negative values ranging in the amplitude \hl{between} -1000 to -200 nT. The scripts used \hl{by} the GMT \hl{modules }were adapted to each processed image for \hl{a }semi-automatic \hl{technical }workflow, allowing the comparison of these data\hl{. The }experimental applications of the \hl{GMT} scripts \hl{validated} the \hl{suitability and effectiveness of the }proposed method in different \hl{tasks of the }multi-source geospatial data processing.  

We have shown that the distribution of the Bangui Magnetic Anomaly correlates with the negative values in gravity and geoid as well as \hl{local }depressions in \hl{a }regional topography \hl{of CAR}. \hl{. The analysis of the correlation between the datasets revealed }the nonlinear relationship\hl{ }between the tectonic evolution \hl{in the past} and \hl{present }geophysical processes, reflected in local variations \hl{of} regional topography \hl{and gravity around the magnetic anomaly}. \hl{Besides the analysis of coherence of the geological and geophysical datasets on CAR, this paper discussed the aspects of the GMT-based mapping with provided scripts as a technical reference. Spatial analysis in Earth sciences relies on data representation in the form of maps. Therefore, advanced mapping techniques support the geological studies by providing new information that reveals hidden structures on the Earth's surface. }

From the cartographic perspective, we proposed \hl{an advanced} framework for mapping several geospatial datasets using \hl{the }GMT scripts to reveal the coherency \hl{between }geophysical \hl{and topographic }data. We\hl{ }evaluated the distribution of values in topographic, gravity and magnetic contrasts over CAR using \hl{several multi-source }high-resolution\hl{ }datasets: EMAG2, WDMAM for magnetic anomalies, GEBCO/SRTM and ETOPO1 for topography, IGPP for gravity grids, EGM2008 for geoid undulations and the USGS data on the geology of CAR. Our experiments show that scripting approach of GMT can significantly reduce the workload of mapping compared to the QGIS while yielding the solution on quality of maps aimed at obtaining the print-quality maps. Our approach is\hl{ }flexible enough to deal with multi-source datasets of various formats and origin as well as rapid plotting of several maps in \hl{the }identical projection and spatial extent for comparability. 

\hl{This paper adds to current knowledge of the CAR region by presenting the integrated approach of mapping the multi-source geophysical, topographic and geological data with repeated structures. In particular, we analysed the coherency of the Bangui magnetic anomaly with gravity and topography data. Similar features in geological and geophysical data were identified and the correlation of the data distribution was discussed to add information in the existing geophysical studies on CAR. The study objectives were implemented and tested on open data from the publicly available geophysical and topographic grids. Analysing the results, we note a comparability in the distribution of geophysical and magnetic data over CAR which is caused by the links between the processes related to the inner geophysical processes in the Earth's crust and the topographic structure.}

%%%%%%%%%%%%%%%%%%%%%%%%%%%%%%%%%%%%%%%%%%
\authorcontributions{Supervision, conceptualization, methodology, software, resources, funding acquisition, and project administration, O.D.; writing---original draft preparation, methodology, software, data curation, visualization, formal analysis, validation, writing---review and editing, and investigation, P.L. All authors have read and agreed to the published version of the manuscript.}

\funding{The publication was funded by the Editorial Office of \emph{Minerals}, Multidisciplinary Digital Publishing Institute (MDPI), by providing 100\% discount for the APC of this manuscript. This project was supported by the Federal Public Planning Service Science Policy or Belgian Science Policy Office, Federal Science Policy - BELSPO (B2/202/P2/SEISMOSTORM).}

\institutionalreview{Not applicable.}

\informedconsent{Not applicable.}

\dataavailability{Not applicable} 

\acknowledgments{The authors thank the three anonymous reviewers for reading and providing useful comments \textcolor{blue}{which improved the initial version of }this manuscript. \textcolor{blue}{The GMT \cite{WesselSmith1991} and QGIS \cite{QGIS_software} software packages were used extensively in the preparation of this paper.}}

\conflictsofinterest{The authors declare no conflict of interest.} 

\abbreviations{Abbreviations}{
The following abbreviations are used in this manuscript:\\

\noindent 
\begin{tabular}{@{}ll}
CAR & Central African Republic \\
CGMW & Commission for the Geological Map of the World \\
EGM2008 & Earth Gravitational Model 2008 \\
EMAG2 & Earth Magnetic Anomaly Grid \\
ETOPO1 & Earth Topography Model with 1 arc minute resolution\\
FFT & Fast Fourier Transform \\
GEBCO & General Bathymetric Chart of the Oceans \\
GMT & Generic Mapping Tools \\
IAGA & International Association of Geomagnetism and Aeronomy \\
IGPP & Institute of Geophysics and Planetary Physics \\
SRTM & Shuttle Radar Topography Mission \\
USGS & United States Geological Survey \\
WDMAM & World Digital Magnetic Anomaly Map \\
\end{tabular}
}

\appendixtitles{no} % Leave argument "no" if all appendix headings stay EMPTY (then no dot is printed after "Appendix A"). If the appendix sections contain a heading then change the argument to "yes".
\appendixstart
\appendix
\section[\appendixname~\thesection]{GMT scripts}
\subsection[\appendixname~\thesubsection]{GMT script for computing the histogram equalisation of CAR's topographic grids}

\begin{lstlisting}[language=bash,caption=GMT script for computing the histogram equalisation of CAR's topographic grids,style=mystyle,label={lst01}]
#!/bin/bash
gmt grdcut ETOPO1_Ice_g_gmt4.grd -R14/28/2/11.5 -Gcf1_relief.nc
ps=HistCAR.ps
gmt makecpt -Crainbow -V -T212/1820 > t.cpt
gmt makecpt -Crainbow -T0/15/1 > c.cpt
# 1 upper left plot
gmt grdhisteq cf1_relief.nc -Gout.nc -C16
gmt grdimage cf1_relief.nc -I+a45+nt1 -Ct.cpt -JM3i -Y6i -K -P \
    -Bpxg4f1a2 -Bpyg4f2a2 -Bsxg2 -Bsyg2 -BWSNe > $ps
echo "24 10 Original" | gmt pstext -Rcf1_relief.nc -J -O -K -F+jBL+f12p -T -Gwhite@10 -Dj0.1i >> $ps
gmt pscoast -R -J -Ia/thinner,blue -Na -N1/thick,white -W0.1p -Df -O -K >> $ps
# 2 upper right plot
gmt grdimage out.nc -Cc.cpt -J -X3.5i -K -O \
    -Bpxg4f2a4 -Bpyg4f2a2 -Bsxg2 -Bsyg2 -BWSNe >> $ps
gmt pscoast -R -J -Ia/thinner,blue -Na -N1/thick,white -W0.1p -Df -O -K >> $ps
echo "24 10 Equalized" | gmt pstext -R -J -O -K -F+jBL+f12p -T -Gwhite@10 -Dj0.1i >> $ps
gmt psscale -Dx0i/-0.4i+jTC+w5i/0.15i+h+e+n -O -K -Ct.cpt -Bg200f10a200 -By+lm >> $ps
# 3 low left plot
gmt grdhisteq cf1_relief.nc -Gout.nc -N
gmt makecpt -Crainbow -T-3/3 > c.cpt
gmt grdimage out.nc -Cc.cpt -J -X-3.5i -Y-3.0i -K -O \
    -Bpxg4f2a4 -Bpyg4f2a2 -Bsxg2 -Bsyg2 -BWSNe >> $ps
gmt grdcontour cf1_relief.nc -R -J -C200 -A500+f7p,26,blue -Wthinnest,blue -O -K >> $ps
echo "23.5 10 Normalized" | gmt pstext -R -J -O -K -F+jBL+f12p -T \
    -UBL/-5p/-4.0c -Gwhite@10 -Dj0.1i >> $ps
gmt pscoast -R -J -Ia/thinner,blue -Na -N1/thick,white -W0.1p -Df -O -K >> $ps
# 4 low right plot
gmt grdhisteq cf1_relief.nc -Gout.nc -Q
gmt makecpt -Crainbow -T0/15 > q.cpt
gmt grdimage out.nc -Cq.cpt -J -X3.5i -K -O \
    -Bpxg4f2a4 -Bpyg4f2a2 -Bsxg2 -Bsyg2 -BWSNe >> $ps
gmt pscoast -R -J -Ia/thinner,blue -Na -N1/thick,white -W0.1p -Df -O -K >> $ps
echo "24 10 Quadratic" | gmt pstext -R -J -O -K -F+jBL+f12p -T -Gwhite@10 -Dj0.1i >> $ps
gmt psscale -Dx0i/-0.4i+w5i/0.15i+h+jTC+e+n -O -K -Cc.cpt -Bg0.5f0.1a1 -By+l"z@-n@-" >> $ps
gmt psscale -Dx0i/-1.0i+w5i/0.15i+h+jTC+e+n -O -K -Cq.cpt -Bg0.5f0.2a1 -By+l"z@-q@-" >> $ps
gmt logo -Dx0.0/-4.5c+o-1.0c/0.2c+w2c -O -K >> $ps
gmt pstext -R0/10/0/15 -JX10/10 -X-8.0 -Y6.5c -N -O \
-F+f14p,Palatino-Roman,black+jLB >> $ps << EOF
0.0 11.5 Central African Republic: histogram equalization of topographic grid:
1.5 10.7 ETOPO1 DEM Global Relief Model 1 arc min resolution
EOF
gmt psconvert HistCAR.ps -A0.5c -E720 -Tj -Z
\end{lstlisting}

\subsection[\appendixname~\thesubsection]{GMT script for mapping the Earth geoid model on CAR}
\begin{lstlisting}[language=bash,caption=GMT script for mapping the Earth geoid model on CAR,style=mystyle,label={lst02}]
#!/bin/bash
gmt grdconvert n00e00/w001001.adf geoid_CF.grd
gdalinfo geoid_CF.grd -stats
gmt makecpt -Cjet.cpt -T-18/18 > colors.cpt
ps=Geoid_CF.ps
gmt grdimage geoid_CF.grd -Ccolors.cpt -R14/28/2/11.5 -JM6.5i -P -Xc -I+a15+ne0.75 -K > $ps
gmt grdcontour geoid_CF.grd -R -J -C0.5 -A1.0+f8p,0,black -Wthinner,dimgray -O -K >> $ps
gmt psbasemap -R -J --MAP_FRAME_AXES=WEsN \
    --FORMAT_GEO_MAP=ddd:mm:ssF \
    --MAP_GRID_CROSS_SIZE_PRIMARY=0.8c \
    --MAP_GRID_CROSS_SIZE_SECONDARY=0.5c \
    -Bpx1f1a1 -Bpyg1f1a1 -Bsxg1 -Bsyg1 \
    -B+t"Earth geoid model on Central African Republic" -O -K >> $ps
gmt psscale -Dg14/1.5+w16.0c/0.15i+h+o0.3/0i+ml+e -R -J -Ccolors.cpt \
    -Bg2f0.2a4+l"Color scale 'jet': dark to light blue, white, yellow and red [C=RGB], -T-18/18]" \
    -I0.2 -By+lm -O -K >> $ps
gmt psbasemap -R -J -Tmg26/3.6+w3.0c+d1.63+t45/10/5+ithin,lightgoldenrod+p0.1p,gold+l+jCM -O -K >> $ps
gmt psbasemap -R -J -Lx14.5c/-2.3c+c50+w300k+l"Mercator projection. Scale (km)"+f \
    -UBL/-0p/-70p -O -K >> $ps
gmt pscoast -R -J -P -Ia/thinnest,blue -Na -N1/thick,white -Wthick,darkslategray -Df -O -K >> $ps
gmt logo -Dx7.0/-3.1+o0.1i/0.1i+w2c -O -K >> $ps
gmt pstext -R0/10/0/15 -JX10/10 -X0.1c -Y4.5c -N -O \
    -F+f10p,25,black+jLB >> $ps << EOF
3.0 10.4 Earth geoid image EGM2008 vertical datum 2.5 min resolution
EOF
gmt psconvert Geoid_CF.ps -A1.0c -E720 -Tj -Z
\end{lstlisting}

\subsection[\appendixname~\thesubsection]{GMT script for mapping the IGPP Earth Free-Air Anomaly on CAR}
\begin{lstlisting}[language=bash,caption=GMT script for mapping the IGPP Earth Free-Air Anomaly on CAR,style=mystyle,label={lst03}]
#!/bin/bash
gmt img2grd grav_27.1.img -R14/28/2/11.5 -Ggrav_CF.grd -T1 -I1 -E -S0.1 -V
gdalinfo grav_CF.grd -stats
gmt makecpt -Chaxby.cpt -T-85/80 > colors.cpt
ps=Grav_CF.ps
gmt grdimage grav_CF.grd -Ccolors.cpt -R14/28/2/11.5 -JM6.0i -P -I+a15+ne0.75 -Xc -K > $ps
gmt grdcontour grav_CF.grd -R -J -C20 -A40 -Wthinner -O -K >> $ps
gmt psbasemap -R -J --MAP_FRAME_AXES=WEsN \
    --FORMAT_GEO_MAP=ddd:mm:ssF \
    -Bpx4f2a2 -Bpyg4f2a2 -Bsxg2 -Bsyg2 \
    -B+t"Free-air gravity anomaly on Central African Republic" -O -K >> $ps
gmt psscale -Dg14/1+w15.0c/0.15i+h+o0.3/0i+ml+e -R -J -Ccolors.cpt \
    -Bg20f1a10+l"Color scale 'haxby' Bill Haxby's color scheme for geoid & gravity [C=RGB] -T-200/200]" \
    -I0.2 -By+l"mGal" -O -K >> $ps
gmt psbasemap -R -J -Tdg15/10.5+w1.5c+f3+l -Lx13.5c/-2.3c+c50+w300k+l"Mercator projection. Scale: km"+f \
    -UBL/-0p/-70p -O -K >> $ps
gmt pscoast -R -J -P -Ia/thinnest,blue -Na -N1/thickest,deeppink1 -Wthin,darkslategray -Df -O -K >> $ps
gmt logo -Dx7.0/-3.1+o0.1i/0.1i+w2c -O -K >> $ps
gmt pstext -R0/10/0/15 -JX10/10 -X0.5c -Y5.0c -N -O \
    -F+f10p,25,black+jLB >> $ps << EOF
2.3 9.2 IGPP Global Earth Free-Air Anomaly, 0.025 mGal resolution
EOF
gmt psconvert Grav_CF.ps -A0.5c -E720 -Tj -Z
\end{lstlisting}

\subsection[\appendixname~\thesubsection]{GMT script for estimating the coherency between geophysical grids on CAR}
\begin{lstlisting}[language=bash,caption=GMT script for estimating the coherency between geophysical grids on CAR,style=mystyle,label={lst04}]
#!/bin/bash
# Extract subsets of img/grd files in geographic format ----------------->
gmt grdcut GEBCO_2019.nc -R14/28/2/11.5 -Gcf_relief.nc
gmt grdconvert n00e00/w001001.adf geoid_CF.grd
gmt img2grd grav_27.1.img -R14/28/2/11.5 -Ggrav_CF.grd -T1 -I1 -E -S0.1 -V
gmt img2grd curv_27.1.img -R14/28/2/11.5 -Ggravvert_CF.grd -T1 -I1 -E -S0.1 -V
gmt grdcut EMAG2_V2.grd -R14/28/2/11.5 -Gcf_mag.nc
gmt grdcut @earth_wdmam_03m -R14/28/2/11.5 -Gcf_mag_wdmam.nc
gmt grdfft cf_mag_wdmam.nc cf_relief.nc -E+wk+n -Na+d+wtmp > cross_spectra1.txt
gmt set FONT_TITLE 12p GMT_FFT kiss
# make cpt for rasters ----------------->
gmt makecpt -Cturbo -T212/1820 > z.cpt # topo
gmt makecpt -Crainbow -T-18/18 > g.cpt # geoid
gmt makecpt -Cjet.cpt -T-40/40 > c.cpt # grav
gmt makecpt -Chaxby.cpt -T-85/80 > d.cpt # vert grav
gmt makecpt -Cwysiwyg.cpt -T-1000/512 > m.cpt # EMAG2
gmt makecpt -Cmag -T-933/375 > n.cpt # WDMAM
ps=Coherency_CF.ps
# map 1 lower left ----------------->
gmt grdimage cf_relief.nc -R14/28/2/11.5 \
    -I+a0+nt1 -JM6.0c -Cz.cpt -P -K -X1.474i -Y1i > $ps
gmt psbasemap -R14/28/2/11.5 -JM6.0c \
    -Ba -BWSne+t"IGPP Global Earth Relief"  -UBL/-5p/-40p \
    --MAP_TITLE_OFFSET=0.1c -O -K >> $ps
# map 2 lower right ----------------->
gmt grdimage geoid_CF.grd -R14/28/2/11.5 \
    -I+a0+nt1 -JM6.0c -Cg.cpt -O -K -X7.5c >> $ps
gmt psbasemap -R14/28/2/11.5 -JM6.0c \
    --MAP_TITLE_OFFSET=0.1c -Ba -BWSne+t"Geoid model EGM2008" -O -K  >> $ps
# map 3 upper left ----------------->
gmt grdimage gravvert_CF.grd -R14/28/2/11.5 -I+a0+nt1 -JM6.0c \
    -Cc.cpt -O -K -X-7.5c -Y6.0c >> $ps
gmt psbasemap -R14/28/2/11.5 -JM6.0c -Ba \
    -BWSne+t"Vertical gravity gradient" \
    --MAP_TITLE_OFFSET=0.1c -O -K >> $ps
# map 4 upper right ----------------->
gmt grdimage grav_CF.grd -R14/28/2/11.5 -I+a0+nt1 -JM6.0c -Cd.cpt -O -K -X7.5c >> $ps
gmt psbasemap -R14/28/2/11.5 -JM6.0c -Ba -BWSne+t"Free-air gravity anomaly" \
    --MAP_TITLE_OFFSET=0.1c -O -K >> $ps
# map 5 upper left EMAG2 ----------------->
gmt grdimage cf_mag.nc -R14/28/2/11.5 -I+a0+nt1 -JM6.0c \
    -Cm.cpt -O -K -X-7.5c -Y6.0c >> $ps
gmt psbasemap -R14/28/2/11.5 -JM6.0c -Ba -BWSne+t"EMAG2 for CAR" \
    --MAP_TITLE_OFFSET=0.1c -O -K >> $ps
# Bangui magnetic anomaly
gmt psbasemap -R -J -D18/5/22/7r -F+pthickest,white -O -K >> $ps
# map 6 upper right ----------------->
gmt grdimage cf_mag_wdmam.nc -R14/28/2/11.5 -I+a0+nt1 -JM6.0c -Cn.cpt -O -K -X7.5c >> $ps
gmt psbasemap -R -J -D18/5/22/7r -F+pthickest,white -O -K >> $ps
gmt psbasemap -R14/28/2/11.5 -JM6.0c -Ba -BWSne+t"WDMAM for CAR" \
    --MAP_TITLE_OFFSET=0.1c -O -K >> $ps
# plot coherency ----------------->
gmt psxy -R3/100/0/1 -JX14.0cl/3.5c -Bx2g3+u"mGal" -Byaf0.04g0.2+l"Coherency@+2@+" \
    -BWSnE+t"Coherency between vertical gradient and free-air gravity anomaly"+gazure \
    --MAP_GRID_PEN_PRIMARY=thinnest,dimgray \
    --MAP_GRID_PEN_SECONDARY=thinner,dimgray \
    -X-7.5cm -Y6.0c cross_spectra.txt -i0,15 -W0.1p -O -K >> $ps
gmt psxy -R -J cross_spectra.txt -i0,15,16 -Sc0.15c -Gred -W0.25p -Ey -O -K >> $ps
gmt logo -Dx7.5c/-2.0c+o-2.0c/-17.5c+w2c -O >> $ps
# Convert to image file by GhostScript
gmt psconvert Coherency_CF.ps -A1.0c -E720 -Tj -Z
# Unix remove files 
rm -f cross.txt *_tmp.nc ?.cpt bbox
\end{lstlisting}

%%%%%%%%%%%%%%%%%%%%%%%%%%%%%%%%%%%%%%%%%%
\begin{adjustwidth}{-\extralength}{0cm}
%\printendnotes[custom] % Un-comment to print a list of endnotes

\reftitle{References}
%\nocite{*}

%=====================================
% References, variant A: external bibliography
%=====================================
\bibliography{Bangui.bib}

%%%%%%%%%%%%%%%%%%%%%%%%%%%%%%%%%%%%%%%%%%
\end{adjustwidth}
\end{document}
